% =====================================================================
%  PAPER V -- predictive safety filtering for LLM-driven UAV agents
%  under adaptive attack. Target: IEEE SaTML (alt: USENIX Security).
%
%  Standalone. Set as Overleaf's main document; pdfLaTeX + BibTeX.
%    pdflatex PaperV_safety_filter_llm_agent && bibtex PaperV_safety_filter_llm_agent \
%      && pdflatex PaperV_safety_filter_llm_agent && pdflatex PaperV_safety_filter_llm_agent
% =====================================================================
\documentclass[11pt,a4paper]{article}
% =====================================================================
%  Shared preamble for the three research proposals.
%
%  Each proposal \input{}s this file and is otherwise standalone, so any
%  one of them can be set as Overleaf's main document and compiled on its
%  own. Nothing here is specific to a single proposal.
% =====================================================================
\usepackage[T1]{fontenc}
\usepackage[utf8]{inputenc}
% Times, matching the parent manuscript. mathptmx is in every TeX Live and on
% Overleaf; lmodern is deliberately not used, since it is absent from minimal
% distributions and this preamble should compile anywhere.
\usepackage{mathptmx}
\usepackage[a4paper,margin=25mm]{geometry}
\usepackage{amsmath,amssymb,amsthm}
\usepackage{booktabs}
\usepackage{tabularx}
\usepackage{xcolor}
\usepackage{titlesec}
\usepackage{enumitem}
\usepackage{microtype}
\usepackage[hidelinks,breaklinks]{hyperref}

% ---- shared style tokens, matching the parent manuscript ------------
\definecolor{netcol}{RGB}{31,78,121}
\definecolor{autocol}{RGB}{18,120,90}
\definecolor{bridgecol}{RGB}{176,96,0}
\definecolor{accent}{RGB}{140,20,20}

\newcommand{\layernet}{\textcolor{netcol}{\textbf{network}}}
\newcommand{\layerauto}{\textcolor{autocol}{\textbf{autonomy}}}
\newcommand{\dataset}[1]{\textsc{#1}}
\newcommand{\code}[1]{\texttt{#1}}
\newcommand{\Real}{\mathbb{R}}
\newcommand{\Norm}[1]{\left\lVert #1\right\rVert}
\newcommand{\MCR}{\mathrm{MCR}}

\newtheorem{claim}{Claim}
\newtheorem{proposition}{Proposition}
\newtheorem{definition}{Definition}
\theoremstyle{remark}
\newtheorem{remark}{Remark}

% ---- section styling -------------------------------------------------
\titleformat{\section}{\large\bfseries\color{netcol}}{\thesection}{0.7em}{}
\titleformat{\subsection}{\normalsize\bfseries}{\thesubsection}{0.6em}{}
\titlespacing*{\section}{0pt}{1.4em}{0.6em}

\setlist[itemize]{leftmargin=1.35em,topsep=0.3em,itemsep=0.25em}
\setlist[enumerate]{leftmargin=1.6em,topsep=0.3em,itemsep=0.25em}

% ---- a boxed "what this rests on" callout ---------------------------
\newenvironment{honestly}
  {\par\medskip\noindent\begin{tabular}{@{}p{2.5mm}@{\hspace{3mm}}p{0.9\linewidth}@{}}
   \textcolor{accent}{\rule{2.5mm}{\baselineskip}} &
   \small\itshape}
  {\end{tabular}\par\medskip}

% ---- the proposal header --------------------------------------------
\newcommand{\proposalheader}[3]{%
  \begin{center}
  {\Large\bfseries #1\par}\medskip
  {\large\itshape #2\par}\medskip
  {\small #3\par}
  \end{center}
  \medskip\hrule\medskip}

\usepackage{array}
\newcolumntype{Y}{>{\raggedright\arraybackslash}X}
\newcolumntype{P}[1]{>{\raggedright\arraybackslash}p{#1}}
\newcommand{\safe}{\mathcal S}
\newcommand{\adm}{\mathcal A_{\mathrm{adm}}}
\newcommand{\prov}[1]{\textsc{\small[#1]}}

\begin{document}

\proposalheader
  {Paper V --- The Untrusted Planner}
  {Keep Flying Whatever the Model Says: Predictive Safety Filtering for
   Language-Model UAV Agents under Adaptive Attack}
  {Research proposal from the RobustUAVs.ai programme $\cdot$
   Roger Nick Anaedevha $\cdot$ 2026-10-01\\
   Target: IEEE SaTML (fallback: USENIX Security) $\cdot$ builds on the parent
   manuscript~\cite{parentpaper} and Paper~III}

\begin{abstract}
\noindent
Language-model agents are beginning to sit in the high-level loop of autonomous
aircraft: translating an operator's instruction into waypoints, summarising
telemetry, choosing a task, proposing a mode change. That places an attackable,
hard-to-secure component upstream of a safety-critical controller, and the
natural question --- can we stop the language model from being fooled? --- has no
general answer. We propose to stop asking it. A \emph{predictive safety filter}
sits between the agent's proposed action and the actuators: it admits the action
only if the certified navigation model (the parent manuscript's) has a
recursively feasible trajectory keeping the aircraft inside its safe envelope
over the planning horizon, and otherwise overrides to a verified fallback. The
security claim is then independent of the language model's internals: for
\emph{any} sequence of agent outputs, adversarial or not, the physical state
stays in the certified safe set. An adaptive attacker who fully controls the
agent's inputs is thereby downgraded from causing a crash to causing a
\emph{stall} --- the filter converts a safety attack into a liveness attack, and
we quantify the residual it cannot prevent. We make three contributions: the
safety-invariance theorem for an untrusted planner; an adaptive-attack
evaluation protocol, following the discipline that a defence must be tested
against an adversary who knows the filter; and the measured trade between the
filter's conservatism and the mission-denial rate an attacker can force, run on
the RobustUAVs.ai simulator with real language-model agents in the loop.
\end{abstract}

% =====================================================================
\section{The problem, and the move}
% =====================================================================
Autonomy stacks increasingly place a language-model agent in the planning
loop~\cite{openai2024agents}: it reads a textual description of the situation ---
operator messages, telemetry summaries, a map, peer reports --- and emits a
high-level action. That channel is attackable. Indirect prompt
injection~\cite{greshake2023injection} lets an adversary who controls any text
the agent ingests --- a spoofed operator note, a tampered telemetry string, a
poisoned map annotation --- steer the agent's output. Securing the language
model against this \emph{in general} is, on current evidence, not something
anyone can promise.

The move this paper makes is to stop trying. Borrowing the logic of the Simplex
architecture~\cite{seto1998simplex,sha2001simplex} and the predictive safety
filter~\cite{wabersich2021psf,hewing2020mpclearning}, we treat the agent as an
\emph{untrusted high-level planner} and interpose a filter that trusts only the
certified navigation model:

\begin{quote}\itshape
The language model may propose anything. The filter executes the proposal only
when a certified safe trajectory exists that begins with it; otherwise it
executes a verified fallback. Physical safety then holds for every possible
agent output, so whether the attacker fooled the model becomes irrelevant to
whether the aircraft stays in its envelope.
\end{quote}

\noindent
This is the same object as Paper~III's projection layer, with two changes that
make it a security paper rather than a control one: the source of the proposed
action is an \emph{adversarially controllable} language model rather than a
trained policy, and the evaluation is against an \emph{adaptive} attacker who
knows the filter.

% =====================================================================
\section{Where MPC and the language model meet}
% =====================================================================
The filter is a model-predictive control
problem~\cite{rawlings2017mpc}. Let the navigation state evolve as
$\dot x=F(x,u)$ ($F$ Lipschitz with constant $L$ over the operating region, as
in the parent manuscript), let $\safe$ be the mission's safe set (corridor
margin $m$, certified radius $R_i(t)\ge R_{\min}$), and let the language model
propose a high-level action $a_t^{\mathrm{LLM}}$ (a waypoint, a task, a mode).

\begin{definition}[Predictive safety filter]
Given $a_t^{\mathrm{LLM}}$, solve
\[
a_t^{\star}=\arg\min_{a}\ \|a-a_t^{\mathrm{LLM}}\|
\]
subject to: a certified trajectory from $x_t$ under $a$ stays in $\safe$ over the
horizon $H$ and ends in a control-invariant terminal set. Execute $a_t^{\star}$
if the problem is feasible; else execute the verified fallback (loiter, or
return-to-launch).
\end{definition}

\noindent
The admissible set $\adm(x_t)$ is exactly the certified set the parent paper's
budget produces --- the filter is where the certificate is \emph{used}, not just
reported. When $a_t^{\mathrm{LLM}}\in\adm$, the agent's action passes untouched;
the filter is transparent in the benign case and binds only when the proposal
would leave the envelope.

% =====================================================================
\section{Three results}
% =====================================================================
\subsection{Safety is invariant to the planner (the headline)}
\begin{proposition}[Safety under an arbitrary planner]
\label{pr:safety}
If the terminal set is control-invariant and the fallback is recursively
feasible from it, then under the filter the physical state satisfies
$x_t\in\safe$ for all $t$, for \emph{every} sequence
$\{a_t^{\mathrm{LLM}}\}$ --- including one chosen by an adversary with full
control of the agent's inputs and full knowledge of the filter.
\end{proposition}
\noindent
This is recursive feasibility of the predictive safety
filter~\cite{wabersich2021psf} carried to an adversarial planner. The content is
not the MPC machinery, which is standard, but the security reading: the
attacker's control of the language model does not appear in the conclusion. The
safety guarantee has been made \emph{independent of the component that cannot be
secured}.

\subsection{The attack surface collapses to liveness}
\begin{proposition}[Safety attack $\to$ liveness attack]
\label{pr:liveness}
Under the filter, an adaptive adversary controlling the agent's inputs cannot
cause a safety violation (Prop.~\ref{pr:safety}); its maximal effect is to force
the filter to override, i.e.\ to induce the fallback. Its power is therefore
bounded by the \emph{induced fallback rate} --- mission denial --- and never
extends to envelope violation.
\end{proposition}
\noindent
This reframing is the paper's point. ``Can the attacker crash the drone?'' (no,
under the filter) is replaced by ``how often can the attacker stall the
mission?'' (a measurable rate), which is both the right security question and one
the platform can answer.

\subsection{The language attack and the physical attack share one budget}
\begin{proposition}[Composition with the perturbation budget]
\label{pr:compose}
A concurrent navigation-degradation attack (jamming, spoofing) shrinks
$\adm(x_t)$ through the same budget $\delta_{\mathrm{tot}}$ of
Paper~IV. The induced fallback rate is therefore monotone in
$\delta_{\mathrm{tot}}$: the language-channel attack and the physical-channel
attack compose, and the filter's conservatism is the common currency.
\end{proposition}
\noindent
This is the bridge to the rest of the programme: Paper~V's attacker and
Paper~IV's attacker draw on one certified budget, exactly as Paper~I's
credit/debit did.

\begin{honestly}
Proposition~\ref{pr:safety} guarantees \emph{safety}, not \emph{liveness}. The
filter keeps the aircraft in its envelope; it does not make the mission succeed,
and a sufficiently aggressive attacker can hold the system in fallback
indefinitely (persistent denial). That is a real limitation and the honest
headline is the pair: \emph{safety is unconditional under the stated model,
liveness is not}. The paper must report the denial rate as prominently as the
zero-violation result, or it oversells.
\end{honestly}

% =====================================================================
\section{Threat model}
% =====================================================================
\begin{itemize}
\item \textbf{Adversary capability.} Full control of any text the agent ingests
  (indirect prompt injection~\cite{greshake2023injection}); full knowledge of
  the filter, the certified model, and the fallback. Optionally, a concurrent
  bounded navigation-degradation attack (Paper~IV's budget).
\item \textbf{Adversary goal, in order of severity.} (i) A physical safety
  violation --- leave the corridor, breach separation. (ii) Failing that,
  mission denial --- maximise the fallback rate. (iii) A \emph{stealthy} denial
  that the operator does not attribute to an attack.
\item \textbf{Trust boundary.} The language model is untrusted. The certified
  navigation model, the filter solver, and the fallback are trusted (and are
  small enough to be verified, which the language model is not). The paper is
  explicit that it secures the \emph{system}, not the model.
\item \textbf{Out of scope.} Attacks on the trusted base (the filter solver, the
  certified model itself) --- these are the assumptions, stated as assumptions.
\end{itemize}

% =====================================================================
\section{Adaptive evaluation (the methodological core)}
% =====================================================================
A safety filter evaluated against a \emph{fixed} attack proves nothing: the
lesson of a decade of adversarial-ML
evaluation~\cite{carlini2017towards,tramer2020adaptive} is that a defence must
be tested against an adversary adaptive to it. We carry that discipline to the
physical setting.

\begin{description}[leftmargin=2.6em,style=nextline]
\item[RQ1] \textbf{Does the filter hold the zero-violation guarantee empirically,
  under an adaptive adversary?} Expected: zero physical violations, matching
  Prop.~\ref{pr:safety}; any violation is a bug in the instantiation (an
  infeasible solve executed, a model-mismatch term exceeded) and is treated as
  one.
\item[RQ2] \textbf{What denial rate can an adaptive adversary force}, as a
  function of its control over the input channel and its knowledge of the
  filter?
\item[RQ3] \textbf{How does denial scale when the language attack composes with
  a navigation-degradation attack} (Prop.~\ref{pr:compose})?
\item[RQ4] \textbf{What does the filter cost when there is no attack?} The
  benign-case intervention rate and mission-time overhead --- the conservatism
  price of the guarantee.
\end{description}

\paragraph{Adaptivity, concretely and defensively.} The adaptive adversary is
constructed to probe the filter's \emph{boundary} --- the inputs most likely to
drive the agent toward the edge of $\adm$ --- because that is what a sound
evaluation of the \emph{defence} requires. The deliverable is the protocol and
the measured denial curve, not an attack toolkit: following the parent
manuscript's handling of its evasion analysis, the mechanism is described at the
level needed to reproduce the evaluation, the result \emph{bounds} what the
attacker achieves, and the practical output is a filter-configuration procedure.

% =====================================================================
\section{What is in hand, what is new}
% =====================================================================
\begin{center}\small
\begin{tabularx}{\linewidth}{@{}P{0.42\linewidth}Y@{}}
\toprule
\textbf{In hand} & \textbf{To build} \\
\midrule
Certificate engine producing $\adm$ (parent; $\gamma=1.625$\,m/s supremum) &
the MPC safety-filter solver over $\adm$ with a verified fallback \\
Fleet simulator (\code{MAX\_FLEET}=12, $T=70$\,s, $\kappa=0.25$) and ten attack
classes & a language-model agent in the planning loop (the platform already
integrates anthropic / openai / gemini / deepseek) \\
Paper~IV's perturbation budget $\delta_{\mathrm{tot}}$ &
the adaptive-attack protocol and the denial-rate measurement \\
Paper~III's projection/feasibility machinery & its lift from discrete allocation
to continuous MPC under an untrusted planner \\
\bottomrule
\end{tabularx}
\end{center}

% =====================================================================
\section{Experimental plan on RobustUAVs.ai}
% =====================================================================
\begin{enumerate}
\item \textbf{Agent in the loop.} Put a language-model planner (the platform's
  existing providers) on the mission-planning task of the fleet simulator:
  it proposes waypoints, task assignments and mode changes from a textual
  situation description.
\item \textbf{Filter.} Implement the predictive safety filter using the
  certificate engine for $\adm$ and a loiter/RTL fallback; log every
  intervention.
\item \textbf{Adaptive attack.} An adversary controls the textual channel feeding
  the agent and adapts to the filter. Measure physical violations (expect zero),
  the induced fallback rate, and stealthy-denial success.
\item \textbf{Composition.} Repeat under a concurrent navigation-degradation
  attack from Paper~IV's budget, to trace the denial surface over
  $\delta_{\mathrm{tot}}$ (Prop.~\ref{pr:compose}).
\item \textbf{Benign cost.} Intervention rate and mission-time overhead with no
  attacker (RQ4).
\item \textbf{Platform page.} A \emph{Safety Filter} page showing the agent's
  proposal, the filtered action, the admissible set, and the live fallback rate
  under attack --- reusing the existing export and provenance machinery.
\end{enumerate}

\begin{honestly}
Model-plant mismatch is the quiet assumption. Prop.~\ref{pr:safety} holds for
the \emph{certified model}; a gap between it and the simulated dynamics must be
carried as an explicit robustness margin in the filter, and on real hardware it
would have to be identified, not assumed. The simulation study can state the
margin it uses; it cannot claim hardware safety, and must not.
\end{honestly}

% =====================================================================
\section{Why this venue, and why it is distinct}
% =====================================================================
\textbf{IEEE SaTML} is the fit: secure and trustworthy machine learning, with
language-model agent security and adaptive-attack evaluation squarely in scope,
and a control-theoretic safety mechanism as the defence. USENIX Security is the
fallback if the systems/benchmark side dominates.

\noindent Distinct from the others in object:
\begin{itemize}
\item \textbf{Paper~III} filters the output of a \emph{trusted} RL policy for
  feasibility; Paper~V filters the output of an \emph{untrusted} language model
  for safety, and is evaluated adversarially.
\item \textbf{Paper~IV} attacks the \emph{sensing} channel; Paper~V attacks the
  \emph{planning} channel, and Prop.~\ref{pr:compose} shows they share a budget.
\item \textbf{Paper~II} would classify Paper~V as a runtime-assurance lift ---
  a per-operation guarantee discharged by a monitor rather than an offline
  certificate --- which is exactly the bridge Paper~II names for objectives no
  static certificate can reach.
\end{itemize}

% =====================================================================
\section{Risks}
% =====================================================================
\begin{itemize}
\item \textbf{``The guarantee is just MPC recursive feasibility.''} True, and
  stated: the novelty is the security reading (safety independent of an
  unsecurable component), the adaptive-attack evaluation, and the composition
  with the sensing budget --- not the MPC theorem.
\item \textbf{Denial might be trivially easy.} If any adaptive attack forces
  near-100\% fallback, the filter guarantees safety but no useful liveness, and
  the honest finding becomes ``safety filtering is necessary but not sufficient;
  the language channel needs its own defence''. That is still a result, and it
  motivates the next paper.
\item \textbf{Adaptive-attack completeness.} We cannot prove we found the
  strongest attack; we report the protocol and the achieved denial, per the
  discipline~\cite{tramer2020adaptive}, and never claim a worst-case denial
  bound we did not prove.
\item \textbf{Dual-use.} The adaptive-attack construction is dual-use. Following
  the parent manuscript, it is presented with its mitigation (the filter) in the
  same work, bounds what the attacker achieves, targets a defensive artifact, and
  is evaluated in simulation against the project's own defensive platform. An
  ethics section is mandatory.
\end{itemize}

% =====================================================================
\section{Page plan and timeline}
% =====================================================================
\begin{center}\small
\begin{tabular}{@{}lr@{\qquad}lr@{}}
\toprule
\textbf{Section} & \textbf{pp} & \textbf{Section} & \textbf{pp} \\
\midrule
Introduction                 & 1.0  & Adaptive evaluation + results & 2.5 \\
Threat model                 & 1.0  & Composition with sensing      & 1.0 \\
Filter and MPC formulation   & 1.5  & Discussion, ethics, limits    & 1.0 \\
Safety-invariance theory     & 1.75 & Related work, conclusion      & 1.0 \\
\midrule
\multicolumn{3}{@{}l}{\textbf{Body (SaTML / USENIX range)}} & \textbf{10.75} \\
\bottomrule
\end{tabular}
\end{center}

\noindent
About five months. The certificate engine, the fleet simulator and the
language-model integration already exist; the new engineering is the MPC filter,
the agent-in-the-loop harness and the adaptive-attack protocol. The safety
theory is standard MPC recursive feasibility, so it can be written early; the
risk and the interest are both empirical --- how much liveness the filter costs,
and how much an adaptive attacker can take away.

\small
\bibliographystyle{plain}
\bibliography{proposals}

\end{document}
